% ============================================================
%  ch3_related_work.tex
%  Chapter 3: Related Work
%  Master's Thesis -- SQL MATCH_RECOGNIZE on Pandas DataFrames
% ============================================================
\chapter{Related Work}
\label{chap:related}

Row pattern recognition was first proposed as a SQL standard feature
by Zemke et al.~\cite{zemke2007pattern} and was officially adopted
in SQL:2016~\cite{ISO2016SQL,iso2016}.
Subsequent editions refined and documented its
semantics~\cite{ISO2021,melton2017sqlrpr,michels2018new}.
The clause has since been studied both as a specification and in
terms of its practical
implementations~\cite{petkovic2022rpr,petkovic2022fraud}.

The \texttt{MATCH\_RECOGNIZE} clause is supported natively by many
leading SQL engines, including
Oracle~\cite{oracle2022,book},
Snowflake~\cite{snowflake2025match_recognize}, and
Trino~\cite{trino2022,trino2021row}.
Other systems are still considering adding this feature, including
PostgreSQL~\cite{howe2020matchrecognize},
MariaDB~\cite{mariadb2025regex}, and
Microsoft SQL Server~\cite{Zhu2023rowpattern}.
To use \texttt{MATCH\_RECOGNIZE} in systems without native support,
users must emulate row-pattern matching using window functions,
user-defined aggregates, or external procedural
layers~\cite{TURP2,TURP,cao2012optimizationanalyticwindowfunctions}---approaches
that are often more complex, less efficient, and harder to maintain.

% ----------------------------------------------------------------
\section{SQL Engines with Native Support}
\label{sec:related-native}
% ----------------------------------------------------------------

Oracle~\cite{oracle2022,book}, Trino~\cite{trino2022,trino2021row},
Snowflake~\cite{snowflake2025match_recognize}, and
Flink~\cite{flink} all provide advanced solutions for SQL-based
pattern matching, but they each come with important practical
challenges that users need to keep in mind.
Although \texttt{MATCH\_RECOGNIZE} was already defined as part of the
SQL:2016 standard, it has only been implemented in a limited number
of DBMSs to date~\cite{petkovic2022rpr}.
Given the construct's capacity to encapsulate the use of numerous
other queries as a simplification, it would be beneficial to extend
its application to more database systems.
The Oracle implementation serves as the primary reference point for
conformance with the SQL standard~\cite{oracle2022,oracle2021dwhsg,book}.

% ................................
\subsection{Oracle}
\label{subsec:related-oracle}
% ................................

Oracle has been a pioneer in SQL \texttt{MATCH\_RECOGNIZE}, having
introduced row pattern recognition in Version~12c.
It provides a comprehensive and robust implementation well-suited
for complex, enterprise-scale analytics, implementing
\texttt{MATCH\_RECOGNIZE} based on Feature~R010.
It executes pattern-matching queries in a highly efficient manner
through internal optimisations, including backtracking algorithms
specifically designed to address complex patterns.
Due to implementation constraints, however, Oracle does not support
Feature~R020---that is, the use of \texttt{MATCH\_RECOGNIZE} within
the \texttt{WINDOW} clause is not permitted.
There are also further restrictions during use: for example, Boolean
expressions are not allowed in the \texttt{MEASURES} clause, and
the \texttt{DISTINCT} keyword cannot be used before aggregates.
Despite its technical capability, Oracle's power comes with high
operational costs, including commercial licensing fees, which makes
it less accessible for many users~\cite{oracle2022,oracle2021dwhsg,book}.

% ................................
\subsection{Trino}
\label{subsec:related-trino}
% ................................

As an open-source project, Trino is defined by its role as a
distributed data analytics tool offering capabilities for exploring
voluminous data sets, where only a limited number of competitors
provide analogous functionality within the same scope.
In contrast to other DBMSs that are predominantly single-node,
Trino is engineered for parallel processing, thereby enabling
management of larger data sets.
The Trino \texttt{MATCH\_RECOGNIZE} implementation is DFA-based,
a distinction that enables deterministic pattern matching in linear
time---a key difference from other engines that may rely on
backtracking.
However, \texttt{MATCH\_RECOGNIZE} can only be used inside the
\texttt{FROM} clause, thereby implementing only Feature~R010.
It is more accessible than Oracle in terms of deployment and cost,
but still requires considerable infrastructure management and may
lack some of Oracle's advanced
optimisations~\cite{trino2022,trino2021row}.

% ................................
\subsection{Snowflake}
\label{subsec:related-snowflake}
% ................................

Snowflake is a prominent cloud-native platform that offers
\texttt{MATCH\_RECOGNIZE} support within its own syntax.
It supports the \texttt{MATCH\_RECOGNIZE} clause both in the
\texttt{FROM} clause and within the \texttt{WINDOW} clause,
covering Features~R010 and R020.
However, in contrast to the SQL standard, Snowflake makes no
syntactic distinction between these two application areas:
a \texttt{WINDOW} clause containing \texttt{MATCH\_RECOGNIZE}
exhibits a strong similarity to its use in \texttt{FROM}, so
Feature~R020 is implemented only semantically rather than through
an explicit syntactic separation.
Snowflake also offers limited support for advanced clauses and has
computational costs that scale with warehouse
size~\cite{snowflake2025match_recognize,hoffa2021snowflake}.

% ................................
\subsection{Apache Flink}
\label{subsec:related-flink}
% ................................

Apache Flink~\cite{vsprem2024building,georgiou2025streaming} is
designed for real-time stream processing, allowing rapid detection
of event patterns in continuous data flows.
While Flink is highly scalable and effective for high-velocity data,
it introduces its own learning curve and is less familiar to users
accustomed to traditional, batch-based SQL systems.
A major limitation shared by all these platforms is the complexity
of integrating SQL-based pattern matching with machine learning
workflows and other analytics
tools~\cite{ML_Workflow,Makrynioti2019sql4mlAD}.
Typically, data must be exported from the database to external
environments for ML modelling, resulting in data duplication,
consistency challenges, and more complicated data
pipelines~\cite{naphade2025evolution}.

% ................................
\subsection{Bringing \texttt{MATCH\_RECOGNIZE} to DuckDB}
\label{subsec:related-duckdb}
% ................................

DuckDB is a relational database management system that offers
several advantages over the DBMSs discussed
above~\cite{raasveldt2019duckdb}.
The most significant advantage is that it does not require a
separate server for operation; instead it can be embedded directly
into the host application, making it well-suited for small projects
and in-process analytics.
DuckDB is open-source, enabling users to contribute new
functionality, and it provides a Python API that permits the
execution of DuckDB queries inside a transpiler for testing and
debugging purposes~\cite{raasveldt2019duckdb}.

Although \texttt{MATCH\_RECOGNIZE} has not yet been implemented in
DuckDB natively, DuckDB's aim for completeness of the SQL interface
means it supports a wide range of other clauses---including the
\texttt{WINDOW} clause and \texttt{WITH RECURSIVE}---which can be
leveraged to transpile \texttt{MATCH\_RECOGNIZE} queries through a
rule-based transformation
approach~\cite{findling2025bringing,mogensen2024compiler,TURP2,TURP}.
This approach filters multiple CTEs, thereby facilitating the
different clauses contained within a \texttt{MATCH\_RECOGNIZE} call.
Any such implementation targeting DuckDB would aim for Feature~R010
of the SQL standard, restricting row pattern matching to the
\texttt{FROM} clause while still supporting a wide range of
aggregate functions inside the \texttt{MEASURES} and
\texttt{DEFINE} clauses.

% ................................
\subsection{Implementation Restrictions Across Systems}
\label{subsec:related-restrictions}
% ................................

Every implementation of the \texttt{MATCH\_RECOGNIZE} clause has
certain restrictions that impose limitations during use.
Common restrictions shared across
systems~\cite{oracle2022,oracle2021dwhsg,trino2022,petkovic2022rpr} include:

\begin{itemize}
  \item The aggregate function \texttt{MATCH\_NUMBER()} is not
    supported in the \texttt{MEASURES} or \texttt{DEFINE} clauses
    in several implementations.
  \item Column access as parameters within aggregate and navigation
    functions is the only permitted form, which limits the
    expressiveness of aggregate functions.
  \item The \texttt{DISTINCT} keyword before aggregates is
    unsupported in Oracle and related systems.
  \item Boolean expressions cannot be used directly in the
    \texttt{MEASURES} clause in Oracle.
\end{itemize}

For many users---especially those with modest analytics needs or who
work independently---installing and managing large-scale systems
like Oracle, Trino, Snowflake, or Flink can be overly complicated
for some users.  These users may only need to perform occasional
pattern detection and would prefer to avoid the overhead of learning
different functions or dealing with varying syntax for similar
queries across platforms~\cite{8509400}.

Table~\ref{tab:engine-comparison} summarises the key characteristics
and feature-level support of these implementations.

\begin{table}[h]
\centering
\caption{Comparison of SQL engines providing \texttt{MATCH\_RECOGNIZE} support}
\label{tab:engine-comparison}
\begin{tabular}{llcccc}
\toprule
\textbf{Engine} & \textbf{Feature Level} & \textbf{Open Source} & \textbf{Deployment} & \textbf{Licence} \\
\midrule
Oracle          & R010 (no R020)  & No  & On-premise / Cloud & Commercial \\
Trino           & R010            & Yes & Distributed        & Apache 2.0 \\
Snowflake       & R010 + R020     & No  & Cloud-native       & Commercial \\
Apache Flink    & CEP library     & Yes & Streaming cluster  & Apache 2.0 \\
DuckDB          & None (planned)  & Yes & In-process         & MIT        \\
Pandas (this work) & R010 (full) & Yes & In-process         & Open-source \\
\bottomrule
\end{tabular}
\end{table}

Additionally, for datasets that fit comfortably in memory, tools
like Pandas offer a faster and more direct alternative, eliminating
the latency of disk-based operations common to larger database
engines~\cite{Evaluation,in-memory_processing}.
Although Oracle, Trino, Snowflake, and Flink each have unique
strengths for SQL pattern matching in large-scale or real-time
scenarios, they introduce extra costs, infrastructure requirements,
and integration issues that can be prohibitive for everyday users.
These limitations highlight the value of Pandas, as it offers
simpler, memory-based tools and unified workflows, where pattern
matching and data analytics can be performed seamlessly and without
unnecessary duplication or complexity.

% ----------------------------------------------------------------
\section{Pandas and the Python Ecosystem Gap}
\label{sec:related-pandas-gap}
% ----------------------------------------------------------------

Pandas is a popular, open-source Python library for data analysis
and manipulation.  Its straightforward installation process and
strong community support make it a good choice for students,
researchers, and professionals who frequently need to analyse data
without deploying complicated systems like Oracle, Trino, or
Snowflake.  For everyday tasks, Pandas offers straightforward
methods for managing data sequences, filtering, grouping, and
performing basic time-series analyses.

However, Pandas has difficulty with advanced pattern detection in
contrast to SQL's \texttt{MATCH\_RECOGNIZE} clause.  It lacks a
straightforward way to express complex patterns; users must create
their own solutions using functions like \texttt{rolling},
\texttt{groupby}, or \texttt{shift}.  These custom approaches can
become complicated and error-prone, especially with advanced
patterns.

Pandas is highly popular due to ease of setup and use, absence of
licensing fees, and smooth integration within the larger Python
libraries.  This makes it great for quickly developing and testing
ideas for typical pattern detection needs.

Despite Python's dominance in data science, no maintained library
implements SQL:2016 \texttt{MATCH\_RECOGNIZE} semantics for Pandas
DataFrames.  Existing approaches require analysts to manually
implement stateful logic using \texttt{DataFrame.apply()} chains,
\texttt{iterrows()}, or external database round-trips---all of which
are error-prone, verbose, and difficult to maintain.  Alternative
DataFrame libraries---Dask, Polars, Modin---similarly lack
sequential pattern primitives.  Even DuckDB, which recently added
partial support for window analytics~\cite{findling2025bringing},
does not yet expose a full \texttt{MATCH\_RECOGNIZE} interface.

% ----------------------------------------------------------------
\section{Contribution of This Work}
\label{sec:related-contribution}
% ----------------------------------------------------------------

This thesis addresses the identified gap with the following
contributions:

\begin{enumerate}
  \item The \textbf{first SQL:2016-compliant
    \texttt{MATCH\_RECOGNIZE} implementation} that runs entirely
    in-process within a Python Pandas workflow with zero additional
    deployment.
  \item \textbf{Complete pattern construct support}, including
    quantifiers, alternation, grouping, anchors, exclusions, and
    \texttt{PERMUTE}, faithfully following the SQL:2016 standard.
  \item A \textbf{production-ready optimisation layer} comprising
    LRU pattern caching, partition-wise memory management, and
    thread-safe concurrent execution.
  \item A \textbf{safe expression evaluation pipeline} that
    prevents code-injection vulnerabilities without sacrificing
    the full expressiveness of user-defined predicates.
  \item \textbf{Empirical validation} on a real-world dataset of
    over 2.2~million rows with 100\% correctness across 25 diverse
    test patterns and throughput reaching up to 12\,918~rows/sec on
    a commodity laptop.
\end{enumerate}

% ................................................................
% END OF CHAPTER
% ................................................................
